% !PRINTED 440
\weberpage{440}
%<<CONT
$\frac{1}{2}(y - \sqrt{-m}\, x)$, and these two quantities must therefore be relatively prime. Consequently equation (12) cannot hold.

The quadratic equation whose two roots are $(lk)$, $(l^{-1}k)$ therefore has % [sic] German print has (18); no antecedent equation (18) occurs in this section here
\gesperrt{only one root} in common with (18), and from this there results the important theorem:

III. \gesperrt{After the adjunction of $\sqrt{D}$, $(lk)$ can be expressed rationally by $(k)$ in the form:}
\weq{13}{(lk) = f_l(k, \sqrt{D}),}
\gesperrt{where $f_l$ is a rational function whose form is determined by the class $l$ alone.}

The change of the sign of $\sqrt{D}$ has the effect that $(lk)$ passes into $(l^{-1}k)$, thus:
\[ f_l(k, \sqrt{D}) = f_{l^{-1}}(k, -\sqrt{D}). \]
If $l$ is two-sided, then by I. the same formula holds, only that $\sqrt{D}$ then does not occur in $f_l$.

If $l'$ is a second class of the same nature as $l$, then one can apply formula (13) to $l'k$ and obtains:
\[ (l\, l' k) = f_l(l' k, \sqrt{D}), \]
which, with the use of
\[ (l' k) = f_{l'}(k, \sqrt{D}) \]
passes into an equation of the form:
\[ (l\, l' k) = f_{l l'}(k, \sqrt{D}) \]
Since on the left side $l$ can be interchanged with $l'$, it follows that
\[ f_{l l'}(k, \sqrt{D}) = f_{l' l}(k, \sqrt{D}), \]
so that $f_{l l'}$ depends only on the composed class $l l'$. In the same way one can derive
\[ (l\, l'\, l'' k) = f_{l l' l''}(k, \sqrt{-m}), \] % The German print switches here from \sqrt{D} to \sqrt{-m}; since D=-m, the notations are equivalent.
and since by \S\,118, II. every arbitrary class $s$ of the discriminant $D$ can be composed out of such classes $l$, formula (13), together with the consequences drawn from it, is no longer bound to the special assumption made about $l$, that the first coefficient be a prime.

\gesperrt{We have then, if $s, k$ denote two arbitrary classes of the discriminant $D$, the formulae:}
\weq{14}{(sk) = f_s(k, \sqrt{D}),}
\weq{15}{(s^{-1}k) = f_s(k, -\sqrt{D}).}
